% =========================================================
% Sequences — Curated Exercise Set (45 Exercises)
% Sources: Abbott, Bartle & Sherbert, Bruckner, Tao, Lebl
% =========================================================

\section*{Sequences — Exercise Review Set}

These exercises develop the fundamental theory of sequences
in $\mathbb{R}$. They focus on convergence, boundedness,
order properties, algebra of limits, and $\varepsilon$–$N$
arguments.

\bigskip

% ---------------------------------------------------------
% Stub Macro
% ---------------------------------------------------------

\newcommand{\seqex}[4]{
\subsection*{Exercise: #1}

\textbf{Source:} #2

\textbf{Technique:} #3

\textbf{Task.} #4

\begin{remark*}[Work Area]
Write your proof or computation here.
\end{remark*}

\bigskip
}

% =========================================================
% ABBOTT
% =========================================================

\section*{Exercises from Abbott}

\seqex
{Abbott 2.2.1}
{Abbott, \textit{Understanding Analysis}}
{tech:unpack-definition}
{Define convergence of a sequence using the $\varepsilon$–$N$ definition.}

\seqex
{Abbott 2.2.2}
{Abbott}
{tech:choose-N}
{Prove that the constant sequence $a_n=c$ converges to $c$.}

\seqex
{Abbott 2.2.3}
{Abbott}
{tech:epsilon-splitting}
{Prove $\lim (1/n)=0$ using the $\varepsilon$–$N$ definition.}

\seqex
{Abbott 2.2.4}
{Abbott}
{tech:unpack-definition}
{Show that limits of sequences are unique.}

\seqex
{Abbott 2.2.5}
{Abbott}
{tech:monotone-bound}
{Prove that every convergent sequence is bounded.}

\seqex
{Abbott 2.2.6}
{Abbott}
{tech:case-splitting}
{Determine whether the sequence $(-1)^n$ converges.}

\seqex
{Abbott 2.2.7}
{Abbott}
{tech:epsilon-splitting}
{Prove $\lim (3+1/n)=3$.}

\seqex
{Abbott 2.2.8}
{Abbott}
{tech:squeeze}
{Show $\lim \sin(1/n)=0$.}

% =========================================================
% BARTLE
% =========================================================

\section*{Exercises from Bartle \& Sherbert}

\seqex
{Bartle 3.1.1}
{Bartle \& Sherbert}
{tech:unpack-definition}
{Define convergence of sequences in $\mathbb{R}$.}

\seqex
{Bartle 3.1.2}
{Bartle \& Sherbert}
{tech:choose-N}
{Prove that $\lim (2/n)=0$.}

\seqex
{Bartle 3.1.3}
{Bartle \& Sherbert}
{tech:epsilon-splitting}
{Show $\lim (n/(n+1))=1$.}

\seqex
{Bartle 3.1.4}
{Bartle \& Sherbert}
{tech:squeeze}
{Show $\lim (\sin n / n)=0$.}

\seqex
{Bartle 3.1.5}
{Bartle \& Sherbert}
{tech:contradiction}
{Prove that a convergent sequence cannot have two different limits.}

\seqex
{Bartle 3.1.6}
{Bartle \& Sherbert}
{tech:monotone-bound}
{Show that every convergent sequence is bounded.}

\seqex
{Bartle 3.1.7}
{Bartle \& Sherbert}
{tech:case-splitting}
{Determine whether the sequence $(-1)^n/n$ converges.}

\seqex
{Bartle 3.1.8}
{Bartle \& Sherbert}
{tech:tail-argument}
{Show that if $a_n\to L$, then $a_{n+1}\to L$.}

% =========================================================
% BRUCKNER
% =========================================================

\section*{Exercises from Bruckner}

\seqex
{Bruckner 2.2.1}
{Bruckner, \textit{Real Analysis}}
{tech:unpack-definition}
{Define convergence of sequences.}

\seqex
{Bruckner 2.2.2}
{Bruckner}
{tech:epsilon-splitting}
{Prove $\lim (1/n^2)=0$.}

\seqex
{Bruckner 2.2.3}
{Bruckner}
{tech:squeeze}
{Show $\lim (\cos n/n)=0$.}

\seqex
{Bruckner 2.2.4}
{Bruckner}
{tech:case-splitting}
{Determine whether $(-1)^n$ converges.}

\seqex
{Bruckner 2.2.5}
{Bruckner}
{tech:unpack-definition}
{Prove uniqueness of limits.}

\seqex
{Bruckner 2.2.6}
{Bruckner}
{tech:monotone-bound}
{Show convergent sequences are bounded.}

\seqex
{Bruckner 2.2.7}
{Bruckner}
{tech:squeeze}
{Show $\lim (\sin(1/n))=0$.}

\seqex
{Bruckner 2.2.8}
{Bruckner}
{tech:choose-N}
{Prove $\lim (5+1/n)=5$.}

% =========================================================
% TAO
% =========================================================

\section*{Exercises from Tao}

\seqex
{Tao 6.1.1}
{Tao, \textit{Analysis I}}
{tech:unpack-definition}
{Define sequence convergence formally.}

\seqex
{Tao 6.1.2}
{Tao}
{tech:epsilon-splitting}
{Prove $\lim (1/n)=0$.}

\seqex
{Tao 6.1.3}
{Tao}
{tech:choose-N}
{Show $\lim (n/(n+1))=1$.}

\seqex
{Tao 6.1.4}
{Tao}
{tech:monotone-bound}
{Show that every convergent sequence is bounded.}

\seqex
{Tao 6.1.5}
{Tao}
{tech:contradiction}
{Prove that limits are unique.}

\seqex
{Tao 6.1.6}
{Tao}
{tech:squeeze}
{Show $\lim (\sin(1/n))=0$.}

\seqex
{Tao 6.1.7}
{Tao}
{tech:case-splitting}
{Determine whether $(-1)^n$ converges.}

\seqex
{Tao 6.1.8}
{Tao}
{tech:tail-argument}
{Show that if $a_n\to L$, then the tail of the sequence also converges to $L$.}

% =========================================================
% LEBL
% =========================================================

\section*{Exercises from Lebl}

\seqex
{Lebl 2.2.1}
{Lebl, \textit{Basic Analysis I}}
{tech:unpack-definition}
{Define convergence of sequences.}

\seqex
{Lebl 2.2.2}
{Lebl}
{tech:choose-N}
{Prove $\lim (2/n)=0$.}

\seqex
{Lebl 2.2.3}
{Lebl}
{tech:epsilon-splitting}
{Show $\lim (3+1/n)=3$.}

\seqex
{Lebl 2.2.4}
{Lebl}
{tech:squeeze}
{Show $\lim (\sin n/n)=0$.}

\seqex
{Lebl 2.2.5}
{Lebl}
{tech:contradiction}
{Show limits are unique.}

\seqex
{Lebl 2.2.6}
{Lebl}
{tech:monotone-bound}
{Show that convergent sequences are bounded.}

\seqex
{Lebl 2.2.7}
{Lebl}
{tech:case-splitting}
{Determine whether $(-1)^n/n$ converges.}

\seqex
{Lebl 2.2.8}
{Lebl}
{tech:tail-argument}
{Show tails of convergent sequences converge to the same limit.}
